\documentclass[11pt,a4paper]{article}
\usepackage[utf8]{inputenc}
\usepackage[T1]{fontenc}
\usepackage{lmodern}
\usepackage[french]{babel}
\usepackage{amsmath,amssymb,mathtools}
\usepackage{icomma}
\usepackage{xcolor}
\usepackage{colortbl}
\usepackage{tikz}
\usepackage[most]{tcolorbox}
\usepackage{enumitem}
\usepackage{array}
\usepackage{booktabs}
\usepackage{fancyhdr}
\usepackage[hidelinks]{hyperref}
\usetikzlibrary{arrows.meta,positioning,calc,patterns,backgrounds,shapes.geometric,decorations.pathreplacing,decorations.markings}
\usepackage[a4paper,margin=2.15cm,headheight=26pt,headsep=12pt]{geometry}
\usepackage{titlesec}

% ---------------- Palette ----------------
\definecolor{primary}{RGB}{0,151,169}
\definecolor{primarydark}{RGB}{0,88,101}
\definecolor{defcol}{RGB}{0,114,178}
\definecolor{thmcol}{RGB}{0,140,120}
\definecolor{formulacol}{RGB}{198,93,9}
\definecolor{excol}{RGB}{0,135,90}
\definecolor{attncol}{RGB}{200,40,55}
\definecolor{methcol}{RGB}{90,90,100}
\definecolor{remarkcol}{RGB}{120,94,200}
\definecolor{barcol}{RGB}{0,151,169}
\definecolor{barcold}{RGB}{0,110,124}
\definecolor{modecol}{RGB}{200,55,55}
\definecolor{meancol}{RGB}{0,90,200}
\definecolor{medcol}{RGB}{160,90,190}
\definecolor{quartcol}{RGB}{0,135,90}
\definecolor{ogivecol}{RGB}{176,74,0}

% ---------------- Boîtes ----------------
\tcbset{boxbase/.style={enhanced,breakable,boxrule=0.9pt,arc=2.5pt,
  left=3mm,right=3mm,top=2.3mm,bottom=2.3mm,
  fonttitle=\bfseries,coltitle=white,toptitle=0.6mm,bottomtitle=0.6mm}}

\newtcolorbox{definition}[1][]{boxbase,colback=defcol!6,colframe=defcol,
  title={\textsc{Définition}\ifx\relax#1\relax\relax\else\textmd{ : \itshape #1}\fi}}
\newtcolorbox{propriete}[1][]{boxbase,colback=thmcol!7,colframe=thmcol,
  title={\textsc{Propriété}\ifx\relax#1\relax\relax\else\textmd{ : \itshape #1}\fi}}
\newtcolorbox{formule}[1][]{boxbase,colback=formulacol!7,colframe=formulacol,
  title={\textsc{Formule}\ifx\relax#1\relax\relax\else\textmd{ : \itshape #1}\fi}}
\newtcolorbox{exemple}[1][]{boxbase,colback=excol!6,colframe=excol,
  title={\textsc{Exemple}\ifx\relax#1\relax\relax\else\textmd{ : \itshape #1}\fi}}
\newtcolorbox{methode}[1][]{boxbase,colback=methcol!8,colframe=methcol,sharp corners,
  title={\textsc{Méthode}\ifx\relax#1\relax\relax\else\textmd{ : \itshape #1}\fi}}
\newtcolorbox{remarque}[1][]{boxbase,colback=remarkcol!7,colframe=remarkcol,
  title={\textsc{Astuce}\ifx\relax#1\relax\relax\else\textmd{ : \itshape #1}\fi}}
\newtcolorbox{attention}[1][]{boxbase,colback=attncol!7,colframe=attncol,
  title={\textsc{Attention}\ifx\relax#1\relax\relax\else\textmd{ : \itshape #1}\fi}}

% Boîte "exercice" et "corrigé" (titre libre passé en argument)
\newtcolorbox{exobox}[1]{boxbase,colback=primary!4,colframe=primary,
  before skip=8pt,after skip=8pt,title={#1}}
\newtcolorbox{corbox}[1]{boxbase,colback=excol!5,colframe=excol!85!black,
  before skip=8pt,after skip=8pt,title={#1}}

% ---------------- En-tête ----------------
\newcommand{\docpart}[1]{\def\thedocpart{#1}}
\docpart{}
\pagestyle{fancy}
\fancyhf{}
\fancyhead[L]{\small\textcolor{primarydark}{\textbf{Fiche 8} — Statistiques descriptives}}
\fancyhead[R]{\small\textcolor{primarydark}{\textsc{\thedocpart}}}
\fancyfoot[C]{\small\thepage}
\renewcommand{\headrulewidth}{0.8pt}
\renewcommand{\headrule}{\hbox to\headwidth{\color{primary}\leaders\hrule height \headrulewidth\hfill}}

\titleformat{\section}{\large\bfseries\color{primarydark}}{\thesection}{0.6em}{}
\titleformat{\subsection}{\normalsize\bfseries\color{primary}}{\thesubsection}{0.6em}{}
\setlist{topsep=2pt,itemsep=1.5pt,parsep=0pt}

% Bandeau de titre du document
\newcommand{\bandeau}[2]{%
\begin{tcolorbox}[boxbase,colback=primary,colframe=primarydark,halign=center,
   before skip=2pt,after skip=12pt]
{\LARGE\bfseries\color{white} #1}\\[3pt]{\normalsize\color{white} #2}
\end{tcolorbox}}

\newcommand{\ecm}{\sigma_m}
\newcommand{\ect}{\sigma}
\newcommand{\med}{M\!e}
\docpart{Difficile — Corrigé — Thème 4}
\begin{document}
\bandeau{Thème 4 — Les indicateurs de dispersion}{Niveau Difficile — corrigé}

\begin{corbox}{Corrigé 1 — Étude de dispersion complète}
\textbf{(a)} Étendue $=170-70=\mathbf{100\ \text{mg/dL}}$.\\[2pt]
\textbf{(b)} $\sum n_i x_i=6(80)+14(100)+18(120)+8(140)+4(160)=480+1400+2160+1120+640=5800$.
\[ m=\frac{5800}{50}=\mathbf{116\ \text{mg/dL}}. \]
\textbf{(c)} $\sum n_i x_i^{\,2}=6(6400)+14(10000)+18(14400)+8(19600)+4(25600)$
\[ =38\,400+140\,000+259\,200+156\,800+102\,400=696\,800. \]
\[ V=\frac{696\,800}{50}-116^2=13\,936-13\,456=\mathbf{480}. \]
\textbf{(d)} $\ect=\sqrt{480}\approx\mathbf{21{,}9\ \text{mg/dL}}$. Interprétation : la glycémie d'un patient
s'écarte, de part et d'autre de la moyenne ($116$), d'environ $21{,}9$\,mg/dL.\\[2pt]
\textbf{(e)} $V=480\geqslant0$ \checkmark. Unités : $V$ en $(\text{mg/dL})^2$, $\ect$ en mg/dL
(même unité que le caractère).
\end{corbox}

\begin{corbox}{Corrigé 2 — Comparer trois écarts types}
\textbf{(a)} Chaque série a pour somme $1+\cdots=36$ (par ex.\ $(1):1+3+5+7+9+11=36$), donc
$m=\dfrac{36}{6}=\mathbf{6}$ pour les trois.\\[2pt]
\textbf{(b)} Écarts au carré autour de $6$, puis $V=\dfrac{1}{6}\sum(x-6)^2$ :
\[
(1):\frac{25+9+1+1+9+25}{6}=\frac{70}{6}\approx11{,}67;\quad
(2):\frac{25+16+1+1+16+25}{6}=\frac{84}{6}=14;
\]
\[
(3):\frac{25+4+1+1+4+25}{6}=\frac{60}{6}=10.
\]
\textbf{(c)} La plus petite variance est celle de la série $\mathbf{(3)}$ ($V=10$), donc le plus petit
écart type. Sans calcul : la série $(3)$ est celle dont les valeurs sont \emph{les plus proches} de
la moyenne $6$ (on y trouve $4,5,7,8$ resserrés autour de $6$).\\[2pt]
\textbf{(d)} La plus grande variance est celle de $(2)$ ($V=14$) : c'est elle qui a le plus grand
écart type (ses valeurs $2$ et $10$ s'éloignent le plus du centre).
\end{corbox}
\end{document}
